% UET Thesis Pandoc LaTeX Template
% Customized for University of Engineering and Technology, VNU

\clearpage
\phantomsection

\addcontentsline{toc}{chapter}{Mở đầu}
\chapter*{Mở đầu}

Hệ mật mã Đường cong Elliptic (ECC) là nền tảng bảo mật của hạ tầng số
hiện đại, từ giao thức TLS, SSH cho đến các hệ thống blockchain
\citep{DiffieHellman1976, Rosing2024}. Tuy nhiên, sự phụ thuộc vào các
bộ tham số sinh sẵn bởi các tổ chức quốc tế đặt ra những thách thức
nghiêm trọng về tự chủ công nghệ và an ninh quốc gia. Đồng thời, hầu hết
các thư viện mật mã truyền thống được viết bằng C/C++ --- một ngôn ngữ
tiềm ẩn nhiều rủi ro về an toàn bộ nhớ.

Khóa luận này đặt mục tiêu thiết kế và cài đặt từ đầu (from scratch) một
hệ thống mật mã đường cong elliptic trên trường nguyên tố cực lớn
1024-bit, sử dụng ngôn ngữ lập trình Rust. Hệ thống bao gồm toàn bộ ngăn
xếp --- từ lớp số nguyên lớn \texttt{U1024}, số học trường hữu hạn
Montgomery, phép toán nhóm trên đường cong, đến các sơ đồ chữ ký số
Schnorr và ECDSA --- không phụ thuộc bất kỳ thư viện mật mã bên ngoài
nào.

\section{Đóng góp của đề
tài}\label{ux111uxf3ng-guxf3p-cux1ee7a-ux111ux1ec1-tuxe0i}

Khóa luận đạt được các đóng góp chính sau:

\begin{enumerate}
\def\labelenumi{\arabic{enumi}.}
\tightlist
\item
  \textbf{Xây dựng hệ thống số học bignum hoàn chỉnh} --- Tự định nghĩa
  kiểu dữ liệu \texttt{U1024} với phép nhân Montgomery, lũy thừa
  hằng-thời-gian, và nghịch đảo Fermat, đảm bảo an toàn bộ nhớ tuyệt đối
  nhờ kiến trúc Rust.
\item
  \textbf{Sinh đường cong elliptic mới bằng thuật toán Cocks-Pinch cải
  tiến} --- Chủ động xây dựng đường cong pairing-friendly (\(k = 18\),
  \(r \approx 2^{512}\), \(p \approx 2^{1024}\)) với cấu trúc
  NTT-friendly trên cả hai trường, đạt mức bảo mật 256-bit và kháng tấn
  công TNFS \citep{CocksPinch2001, BarbulescuDuquesne2019}.
\item
  \textbf{Cài đặt ba sơ đồ chữ ký số} --- Schnorr, ECDSA và chữ ký tổng
  hợp BLS, với bộ kiểm thử 36 test bao phủ từ số học cấp thấp đến mô
  phỏng tấn công MOV và Anomalous.
\end{enumerate}

\section{Bố cục của khóa
luận}\label{bux1ed1-cux1ee5c-cux1ee7a-khuxf3a-luux1eadn}

\begin{itemize}
\tightlist
\item
  \textbf{Chương 1. Tổng quan và Giới thiệu bài toán}: Bối cảnh lịch sử
  của ECC, thực trạng các tiêu chuẩn hiện hành, vấn đề tự chủ mật mã,
  vai trò của ngôn ngữ lập trình, và phát biểu bài toán.
\item
  \textbf{Chương 2. Cơ sở toán học}: Trường hữu hạn, đường cong elliptic
  Short Weierstrass, phép cộng điểm, nhân vô hướng, và phép nhân
  Montgomery.
\item
  \textbf{Chương 3. Phương pháp xây dựng đường cong elliptic}: Phương
  pháp Nhân phức (CM), thuật toán Cocks-Pinch truyền thống và cải tiến
  NTT-friendly.
\item
  \textbf{Chương 4. Sơ đồ chữ ký số}: Cấu trúc cặp khóa, Schnorr, ECDSA,
  và chữ ký tổng hợp BLS.
\item
  \textbf{Chương 5. Cài đặt và Đánh giá}: Benchmark hiệu năng, phân tích
  an toàn trước 5 loại tấn công, và định hướng phát triển.
\end{itemize}

\FloatBarrier
